% frontmatter/english/como-usar.tex -- instrucciones para el adulto, en
% inglés como todo el cuaderno (ver notes/05-english.md).
\section*{How to use this book}

This book is for a child who can already read in Spanish and is learning
English at school. It has one page for every weekday of the year -- 260
pages, holidays too -- and every page has the same three parts:

\begin{itemize}[leftmargin=6mm, itemsep=1mm]
\item \textbf{Date}, \textbf{Read} and \textbf{Notes}, for the grown-up.
  Tick how the reading went: alone, with help, or with difficulty.
\item \textbf{The text of the day}, in the green box: about 30 words in
  autumn and about 70 in summer. The sentences are short, but many of
  them have two parts, joined by \textit{because}, \textit{when},
  \textit{who}, \textit{that}, \textit{but} or \textit{so} -- that is
  what this book practises. The title of the box says how to read it.
\item \textbf{An activity about the text}, in the box below. Almost every
  day there is something to draw, colour, tick, circle or match, and to
  do it the child has to go back to the text: how many, what colour,
  where, why, what happened first. The golden rule is that \textbf{the
  answer is always in the text}. If the child can't find it, don't give
  it away: ask ``Where does it say that?'' and read the text again
  together.
\end{itemize}

\textbf{The activities.} \textbf{Read and draw}: draw exactly what the
text says -- a beautiful drawing is not the point, a true one is -- and
then check it with the questions at the bottom of the box.
\textbf{Read and colour}: colour a picture the way the text says.
\textbf{Where is it?}: draw things in the right place (\textit{in},
\textit{on}, \textit{under}, \textit{next to}...). \textbf{Yes or no?},
\textbf{Read and circle} and \textbf{Read and match}: short questions
about the details. \textbf{Match the halves}: join the two parts of a
long sentence from the text. \textbf{What happened first?}: put the
story in order. \textbf{The missing words}: find a sentence in the text
and write the word that is missing. \textbf{Word hunt}, \textbf{Riddle},
\textbf{Draw the story} (three pictures: the beginning, the middle and the
end of the week), \textbf{Look back} and \textbf{Your turn}. The next page
shows the words and the pictures of every instruction: read it together
on the first day.

\textbf{The story.} The pages tell a year in the life of
\textbf{Lucía}, her little brother \textbf{Dani}, their dog
\textbf{Toby} and their family -- and of their new neighbours, who come
from London and only speak English. Every week is a new chapter.

\textbf{Ten minutes a day.} In autumn, read the text together first and
then let the child read it alone. When a word is new, explain it with a
gesture or a drawing rather than a translation, if you can. Reading a
page twice is always fine; so is doing the activity the next morning.
There is a medal at the end of each part of the year, a diploma at the
end of the book, and an answer key at the back -- for the grown-up, to
check, or to give a small clue.
\newpage
